\documentclass[11pt,a4paper,openany]{ctexbook}
\usepackage[margin=2.2cm]{geometry}
\usepackage{amsmath,amssymb,mathtools}
\usepackage{booktabs,longtable,tabularx,array}
\usepackage[table]{xcolor}
\usepackage[most]{tcolorbox}
\usepackage{enumitem}
\usepackage{fancyhdr}
\usepackage{hyperref}
\usepackage{multicol}
\usepackage{graphicx}

\definecolor{mainteal}{HTML}{0F8B8D}
\definecolor{darkink}{HTML}{15202B}
\definecolor{muted}{HTML}{5E6B75}
\definecolor{softteal}{HTML}{EAF7F7}
\definecolor{softamber}{HTML}{FFF5DD}
\definecolor{softred}{HTML}{FFF0F0}
\definecolor{linegray}{HTML}{DDE4EA}

\hypersetup{
  colorlinks=true,
  linkcolor=mainteal,
  urlcolor=mainteal,
  citecolor=mainteal,
  pdftitle={Derivatives Pricing Review Notes},
  pdfauthor={Course material summary}
}

\pagestyle{fancy}
\fancyhf{}
\lhead{衍生工具定价 / Derivatives Pricing}
\rhead{\thepage}
\renewcommand{\headrulewidth}{0.4pt}
\setlength{\headheight}{14pt}

\setlength{\parindent}{2em}
\setlength{\parskip}{0.35em}
\setlist[itemize]{leftmargin=1.6em,itemsep=0.25em,topsep=0.25em}
\setlist[enumerate]{leftmargin=1.8em,itemsep=0.25em,topsep=0.25em}
\renewcommand{\arraystretch}{1.2}

\newcommand{\term}[2]{\textbf{#1}（\emph{#2}）}
\newcommand{\important}[1]{\textcolor{mainteal}{\textbf{#1}}}
\newcommand{\source}[1]{\par\smallskip\noindent{\small\textcolor{muted}{Source: #1}}\par}

\newtcolorbox{intuitionbox}[1]{
  breakable,
  colback=softteal,
  colframe=mainteal!65!black,
  title=\textbf{#1},
  arc=2mm,
  boxrule=0.7pt
}

\newtcolorbox{formulabox}[1]{
  breakable,
  colback=softamber,
  colframe=orange!70!black,
  title=\textbf{#1},
  arc=2mm,
  boxrule=0.7pt
}

\newtcolorbox{warningbox}[1]{
  breakable,
  colback=softred,
  colframe=red!55!black,
  title=\textbf{#1},
  arc=2mm,
  boxrule=0.7pt
}

\newtcolorbox{workflowbox}[1]{
  breakable,
  colback=white,
  colframe=mainteal!55!black,
  title=\textbf{#1},
  arc=2mm,
  boxrule=0.6pt
}

\newtcolorbox{examplebox}[1]{
  breakable,
  colback=softamber!70!white,
  colframe=orange!65!black,
  title=\textbf{#1},
  arc=2mm,
  boxrule=0.6pt
}

\begin{document}

\frontmatter
\begin{titlepage}
  \centering
  \vspace*{2cm}
  {\Huge\bfseries 衍生工具定价复习讲义\par}
  \vspace{0.5cm}
  {\Large Derivatives Pricing Review Notes\par}
  \vspace{1.2cm}
  {\large 中文主讲 + English key terms + assignment based practice\par}
  \vfill
  \begin{tabular}{ll}
    Course material & Lecture slides, syllabus, assignments, official solutions \\
    Style & Beginner friendly, formula focused, exam oriented \\
    Build & XeLaTeX / ctexbook \\
  \end{tabular}
  \vfill
  {\large Generated from local course files under Course Material\par}
\end{titlepage}

\tableofcontents

\chapter*{How To Use This Review / 怎么用这份讲义}
\addcontentsline{toc}{chapter}{How To Use This Review / 怎么用这份讲义}

这份讲义不是把课件逐页翻译，而是把所有 lecture、assignment 和 solution 读完后，按考试复习逻辑重新组织。每一章都遵循同一个顺序：

\begin{enumerate}
  \item \textbf{先讲直觉}：先回答“这个工具解决什么问题”，避免一上来就背公式。
  \item \textbf{再讲公式}：每个公式都说明变量含义、适用条件和容易错的地方。
  \item \textbf{再做题}：题目主要参考 assignments，assignment 1-3 用官方 solution 核对，assignment 4 根据 Lecture 12 的 binomial tree 方法推导。
  \item \textbf{最后记英文表达}：考试题是英文时，最重要的是认出关键词，例如 forward price, hedge ratio, zero rate, swap rate, put-call parity, risk-neutral probability。
\end{enumerate}

\begin{warningbox}{学习顺序建议}
如果你是小白，先按这个顺序复习：Futures mechanics $\rightarrow$ Hedging $\rightarrow$ Interest rates $\rightarrow$ Forward/Futures pricing $\rightarrow$ Interest rate futures $\rightarrow$ Swaps $\rightarrow$ Options $\rightarrow$ Binomial trees。不要先背 option strategy 图形，否则容易只记形状不理解现金流。
\end{warningbox}

\chapter*{Source Coverage / 材料覆盖清单}
\addcontentsline{toc}{chapter}{Source Coverage / 材料覆盖清单}

\begin{longtable}{p{0.38\textwidth}p{0.18\textwidth}p{0.34\textwidth}}
\toprule
\textbf{Material} & \textbf{Scale} & \textbf{Used for} \\
\midrule
Derivatives Pricing Lecture 1 & 59 slides & course overview, derivative definition, market participants, exchange vs OTC, basic option/forward examples \\
Lecture 2 Futures Markets & 66 slides & futures contract mechanics, margin, clearing house, basis, hedge ratio, stock index futures hedge \\
Lecture 3 Interest Rates part 1 & 43 slides & rate types, LIBOR/SOFR context, compounding, zero rates, bond pricing, bootstrapping \\
Lecture 4 Interest Rates part 2 & 42 slides & forward rates, FRA payoff and valuation, duration and convexity basics \\
Lecture 5 Forward and Futures Prices & 47 slides & no-arbitrage pricing, investment vs consumption assets, income, dividends, storage cost, convenience yield \\
Lecture 6 Interest Rate Futures part 1 & 37 slides & day count, Treasury bond quotation, accrued interest, conversion factor, cheapest-to-deliver \\
Lecture 7 Interest Rate Futures part 2 & 36 slides & Eurodollar futures, duration hedge, convexity adjustment intuition, China rate futures context \\
Lecture 8 Swaps & 48 slides & plain vanilla swaps, comparative advantage, valuation as bonds or FRAs, currency swaps \\
Lecture 9 Securitization & 10 slides & ABS, tranches, waterfall, ABS CDO loss allocation \\
Lecture 10 Option Properties & 49 slides & option payoffs, bounds, early exercise, put-call parity \\
Lecture 11 Trading Strategies & 37 slides & covered calls, protective puts, spreads, straddles, strips, straps, strangles \\
Lecture 12 Binomial Trees & 34 slides & one-step and two-step binomial pricing, risk-neutral valuation, American option early exercise \\
Supplementary Oil Futures Case & 12 slides & WTI/Brent/SC, term structure, geopolitical shock, backwardation/contango interpretation \\
Syllabus / 教学实施方案 & 8 pages & course logistics, textbook, evaluation, reading map \\
Assignments 1-4 & 4 files & practice question bank and final review problems \\
Assignment 1-3 solutions & 3 PDFs & official answer checks for solved practice \\
\bottomrule
\end{longtable}

\mainmatter

\chapter{Big Picture: Derivatives Pricing / 总体框架}
\source{Lecture 1; syllabus}

\section{什么是衍生品 / What Is a Derivative}
\term{衍生品}{derivative} 是一种金融合约，它的价值来自某个 \term{标的资产}{underlying asset}。标的可以是股票、外汇、利率、债券、商品、信用风险，甚至天气指数。课程里的核心工具包括 \term{forward}{远期}, \term{futures}{期货}, \term{swap}{互换}, \term{option}{期权}, 以及基于这些工具组合出来的结构化产品。

\begin{intuitionbox}{小白直觉}
衍生品不是“凭空赚钱”的工具，而是把未来不确定的价格、利率或现金流变成今天可以交易的合约。pricing 的核心问题是：如果市场没有免费午餐（no arbitrage），这个合约今天应该值多少钱？
\end{intuitionbox}

\section{三类使用者 / Three Types of Users}
\begin{itemize}
  \item \term{Hedgers}{套期保值者}：目标不是赌方向，而是减少风险。例如航空公司担心油价上涨，可以买入燃油相关 futures。
  \item \term{Speculators}{投机者}：主动承担方向风险，希望从价格变化中获利。
  \item \term{Arbitrageurs}{套利者}：利用价格关系错位做无风险或近似无风险交易。无套利定价的公式大多来自这个思想。
\end{itemize}

\section{Exchange vs OTC}
{\small
\begin{tabular}{@{}>{\raggedright\arraybackslash}p{0.15\textwidth}>{\raggedright\arraybackslash}p{0.37\textwidth}>{\raggedright\arraybackslash}p{0.37\textwidth}@{}}
\toprule
 & \textbf{Exchange traded} & \textbf{OTC} \\
\midrule
中文 & 交易所交易 & 场外交易 \\
标准化 & 高，合约规格固定 & 低，可定制 \\
信用风险 & 由 clearing house 降低 & 取决于 counterparty \\
典型例子 & futures, exchange-traded options & forwards, swaps, customized options \\
\bottomrule
\end{tabular}
}

\section{课程主线}
这门课的主线可以压缩为四句话：
\begin{enumerate}
  \item Forward/Futures pricing 用 \important{cash-and-carry arbitrage} 建立现货和远期价格的关系。
  \item Interest rate 部分提供 discounting, zero curve, forward rate, duration 等定价基础。
  \item Swaps 可以拆成债券组合或 FRA 组合。
  \item Options 不能只用线性复制，必须处理非线性 payoff；put-call parity 和 binomial trees 是入门核心。
\end{enumerate}

\chapter{Futures Markets and Hedging / 期货市场与套期保值}
\source{Lecture 2; Assignment 1}

\section{Futures Contract Mechanics}
\term{Futures contract}{期货合约} 是交易所标准化合约。合约会规定：
\begin{itemize}
  \item underlying asset：交割什么资产；
  \item contract size：一张合约对应多少数量；
  \item delivery month/place：何时何地交割；
  \item price quotation：价格如何报价；
  \item daily settlement：每天按结算价把盈亏打进 margin account。
\end{itemize}

\begin{intuitionbox}{为什么 futures 要 daily settlement}
Forward 的盈亏通常到期一次性结算；futures 每天结算。这样 clearing house 可以把信用风险切成每天的小块。如果价格朝你不利方向变动，margin account 下降到 maintenance margin 以下，就会触发 margin call。
\end{intuitionbox}

\section{Margin Account / 保证金}
\begin{itemize}
  \item \term{Initial margin}{初始保证金}：开仓时先放进去的钱。
  \item \term{Maintenance margin}{维持保证金}：账户余额不能低于这个数。
  \item \term{Margin call}{追加保证金通知}：余额低于 maintenance margin 时，要补回 initial margin。
\end{itemize}

\begin{formulabox}{Margin call threshold}
对于 long futures：
\[
\text{margin balance}=\text{initial margin}+(\Delta F)\times \text{contract size}\times N.
\]
如果价格下降导致 balance $<$ maintenance margin，就会 margin call。对于 short futures，价格上涨导致亏损。
\end{formulabox}

\section{Basis and Convergence}
\term{Basis}{基差} 通常定义为
\[
\text{basis}=S-F,
\]
其中 $S$ 是 spot price，$F$ 是 futures price。到交割期附近，futures price 和 spot price 必须收敛，否则会出现交割套利。Lecture 2 的 convergence 图就是在说明这一点。

\section{Hedging with Futures}
最小方差套保比率：
\[
h^\ast=\rho\frac{\sigma_S}{\sigma_F}.
\]
这里 $\rho$ 是现货价格变化与期货价格变化的相关系数，$\sigma_S$ 是被套保资产价格变化标准差，$\sigma_F$ 是期货价格变化标准差。

要交易的合约数量：
\[
N^\ast=\frac{h^\ast Q_A}{Q_F},
\]
其中 $Q_A$ 是需要套保的现货数量，$Q_F$ 是一张 futures contract 的数量。

\begin{warningbox}{常见错误}
很多人会把 hedge ratio 的方向和合约数量混在一起。先判断方向：担心价格上涨就需要 long futures，担心价格下跌就需要 short futures；再用 $N^\ast$ 算张数。
\end{warningbox}

\section{Changing Portfolio Beta}
用 stock index futures 改变组合 beta：
\[
N^\ast=(\beta_T-\beta_P)\frac{V_A}{F_0\times \text{multiplier}}.
\]
如果 $N^\ast>0$，买入 futures；如果 $N^\ast<0$，卖出 futures。

\chapter{Interest Rates / 利率、零息利率与远期利率}
\source{Lecture 3; Lecture 4; Assignment 1}

\section{Interest Rate Is a Price of Time}
利率是时间的价格。今天的 1 元和一年后的 1 元不一样，因为今天的钱可以投资。衍生品定价中最常用的是 \term{risk-free rate}{无风险利率}，它用来 discount 确定性现金流。

\section{Compounding / 复利方式}
如果年利率为 $R$，每年复利 $m$ 次，$T$ 年后的金额为
\[
A=P\left(1+\frac{R}{m}\right)^{mT}.
\]
连续复利（continuous compounding）为
\[
A=Pe^{RT}.
\]

离散复利和连续复利转换：
\[
R_c=m\ln\left(1+\frac{R_m}{m}\right),\qquad
R_m=m\left(e^{R_c/m}-1\right).
\]

\begin{intuitionbox}{为什么课程喜欢 continuous compounding}
连续复利让很多公式变短。例如 discount factor 是 $e^{-rT}$，forward price 是 $S_0e^{rT}$。考试给出 continuously compounded rate 时，通常就是暗示你用指数形式。
\end{intuitionbox}

\section{Zero Rate and Discount Factor}
\term{Zero rate}{零息利率} 是从今天到某个期限的一次性投资收益率。若 $R(T)$ 是连续复利 zero rate，discount factor 为
\[
P(0,T)=e^{-R(T)T}.
\]
有现金流 $C_i$ 在 $t_i$ 支付，则价格为
\[
\text{Price}=\sum_i C_iP(0,t_i).
\]

\section{Bootstrapping Zero Curve}
Bootstrapping 是用短期限债券价格一步一步推出更长期 zero rates。思路如下：
\begin{enumerate}
  \item 先用最短期 zero-coupon bond 直接得到第一个 discount factor。
  \item 再用有 coupon 的债券价格，扣掉已经知道的早期 coupon 现值。
  \item 剩下的部分对应最后一期 principal + coupon，由此解出新的 discount factor。
\end{enumerate}

\section{Forward Rate / 远期利率}
如果 $R_1$ 是到 $T_1$ 的 zero rate，$R_2$ 是到 $T_2$ 的 zero rate，均为连续复利，则 $T_1$ 到 $T_2$ 的 forward rate 为
\[
R_F=\frac{R_2T_2-R_1T_1}{T_2-T_1}.
\]

\begin{intuitionbox}{Forward rate 的直觉}
把钱从 0 存到 $T_2$，应该等价于先从 0 存到 $T_1$，再从 $T_1$ 存到 $T_2$。否则就可以借短贷长或借长贷短做套利。
\end{intuitionbox}

\section{Forward Rate Agreement / FRA}
\term{FRA}{forward rate agreement} 是约定未来某段借贷期适用利率的合约。若约定利率为 $R_K$，未来市场利率为 $R_M$，本金 $L$，借贷期长度为 $\tau=T_2-T_1$，则在 $T_2$ 的利差现金流近似为
\[
L(R_K-R_M)\tau
\]
或相反方向，取决于是 lender 还是 borrower。市场惯例通常在 $T_1$ 结算，所以要从 $T_2$ 折现回 $T_1$。

\section{Duration / 久期}
\term{Duration}{久期} 衡量债券价格对收益率变化的敏感性。Modified duration 近似关系：
\[
\frac{\Delta B}{B}\approx-D\Delta y.
\]
Lecture 6-7 在利率期货套保里使用 duration matching：
\[
N^\ast=\frac{P D_P}{F D_F}.
\]

\chapter{Forward and Futures Pricing / 远期与期货定价}
\source{Lecture 5; Assignment 2; supplementary oil futures case}

\section{No Arbitrage Pricing}
无套利定价的基本句式是：构造两个未来现金流完全一样的组合，它们今天价格必须一样。Forward/Futures pricing 最常见的复制组合是：
\begin{itemize}
  \item 买入现货并融资持有到期；
  \item 或签订 forward 到期买入资产。
\end{itemize}

\section{Investment Asset with No Income}
对于不支付收益、无 storage cost 的 investment asset：
\[
F_0=S_0e^{rT}.
\]
如果 $F_0>S_0e^{rT}$，做 cash-and-carry：借钱买现货，同时 short forward。若 $F_0<S_0e^{rT}$，做 reverse cash-and-carry：short spot，把钱投资，同时 long forward。

\section{Known Cash Income}
如果持有资产期间有确定现金收益，其现值为 $I$：
\[
F_0=(S_0-I)e^{rT}.
\]
收益降低 forward price，因为买现货的人在持有期会收到这部分收益。

\section{Known Dividend Yield}
如果资产提供连续收益率 $q$：
\[
F_0=S_0e^{(r-q)T}.
\]
Stock index futures 常用这个公式，其中 $q$ 是 dividend yield。

\section{Foreign Currency}
外汇可以看作支付 foreign risk-free rate $r_f$ 的资产：
\[
F_0=S_0e^{(r_d-r_f)T}.
\]
$r_d$ 是 domestic interest rate，$r_f$ 是 foreign interest rate。

\section{Commodity Futures: Storage Cost and Convenience Yield}
商品分为 \term{investment assets}{投资资产} 和 \term{consumption assets}{消费资产}。黄金更像 investment asset，原油、小麦等更像 consumption asset。若 storage cost 的现值为 $U$，无 convenience yield 时：
\[
F_0=(S_0+U)e^{rT}.
\]
若用连续 storage cost rate $u$ 和 convenience yield $y$：
\[
F_0=S_0e^{(r+u-y)T}.
\]

\begin{intuitionbox}{Convenience yield 的直觉}
持有实物商品有时有额外好处，例如炼油厂持有原油可以保证生产不断。这个“持有实物的便利”会让现货更有价值，因此降低 futures price，相当于公式里的 $-y$。
\end{intuitionbox}

\section{Contango and Backwardation}
\begin{itemize}
  \item \term{Contango}{正向市场}：远期价格高于近月价格，常见原因是 financing/storage cost 较高。
  \item \term{Backwardation}{反向市场}：近月价格高于远期价格，常见原因是短期供给紧张、convenience yield 高。
\end{itemize}
补充油价课件用 WTI、Brent、SC 的 futures curve 说明：地缘冲击会先影响近月合约，因为市场最担心的是即时供给和运输瓶颈，而不是很远未来的均衡价格。

\chapter{Interest Rate Futures / 利率期货}
\source{Lecture 6; Lecture 7; Assignment 2}

\section{Day Count Convention}
\term{Day count convention}{日计数惯例} 决定利息如何按天累计。
\begin{itemize}
  \item Treasury bonds: often Actual/Actual。
  \item Corporate bonds: often 30/360。
  \item Money market instruments: often Actual/360。
\end{itemize}
题目里如果给出 quoted price、coupon date、settlement date，必须先判断 accrued interest。

\section{Treasury Bond Quotation}
美国 Treasury bond 常用 32 分之一报价。例如 102-07 表示
\[
102+\frac{7}{32}=102.21875.
\]
\term{Cash price}{现金价格} 等于 quoted price 加 accrued interest：
\[
\text{Cash price}=\text{Quoted price}+\text{Accrued interest}.
\]

\section{Conversion Factor and Cheapest-to-Deliver}
Treasury bond futures 允许交割多种债券。为了公平，每只债券有 \term{conversion factor}{转换因子}。交割时收到：
\[
\text{Cash received}=\text{Futures price}\times \text{conversion factor}+\text{accrued interest}.
\]
卖方会选择 \term{cheapest-to-deliver}{最便宜可交割债券}，即交割成本最低的债券。

\section{Eurodollar Futures}
Eurodollar futures 的报价：
\[
Q=100-R,
\]
其中 $R$ 是三个月 LIBOR 的年化百分比。若 quote 为 98.40，则 implied rate 为 $1.60\%$。

\begin{warningbox}{Eurodollar hedge direction}
借款人担心未来利率上升。利率上升会使 Eurodollar futures price 下降，所以借款人应该 short Eurodollar futures，价格下跌时 futures 盈利可以抵消借款成本上升。
\end{warningbox}

\section{Duration Based Hedge}
债券组合用 Treasury bond futures 免疫利率风险：
\[
N^\ast=\frac{P D_P}{F D_F}.
\]
这里 $P$ 是组合价值，$D_P$ 是组合 duration，$F$ 是 futures contract 对应的交割价值，$D_F$ 是 CTD bond duration。若担心利率上升导致债券价格下降，通常 short futures。

\chapter{Swaps / 互换}
\source{Lecture 8; Assignment 3}

\section{Plain Vanilla Interest Rate Swap}
\term{Plain vanilla interest rate swap}{普通利率互换} 是一方支付固定利率、收取浮动利率；另一方相反。本金通常不交换，只用 notional principal 计算利息。

\begin{intuitionbox}{Swap 的直觉}
Swap 可以看成一串 FRA，也可以看成一只 fixed-rate bond 和一只 floating-rate bond 的组合。哪个角度更方便，取决于题目给了 swap rate 还是 zero curve。
\end{intuitionbox}

\section{Using Swaps to Transform Liabilities}
如果公司本来借的是 floating rate loan，又进入 swap 支付 fixed、收 LIBOR，那么：
\[
\text{floating loan payment}+\text{swap net payment}
\]
可以合成 fixed rate liability。反过来，fixed loan 加 receive fixed/pay floating swap 可以合成 floating liability。

\section{Comparative Advantage}
比较优势题的套路：
\begin{enumerate}
  \item 分别比较两家公司在 fixed market 和 floating market 的借款差距。
  \item 差距更大的市场是有比较优势的市场。
  \item 总潜在收益等于两个市场差距之差。
  \item 通过 swap 和金融中介分配收益。
\end{enumerate}

\section{Swap Valuation as Bonds}
对“pay fixed, receive floating”的一方：
\[
V_{\text{swap}}=B_{\text{floating}}-B_{\text{fixed}}.
\]
对“receive fixed, pay floating”的一方，符号相反。

\section{Swap Rate and Zero Curve}
Par swap rate 是让 swap 初始价值为 0 的 fixed rate。若支付频率为一年一次，notional 为 100，则 $n$ 年 par swap 可视为 coupon bond at par：
\[
100=\sum_{i=1}^{n-1}\frac{R_{\text{swap}}\times100}{(1+R_i)^i}
\,+\,\frac{100+R_{\text{swap}}\times100}{(1+R_n)^n}.
\]
Assignment 3 第 1 题就是用 2 年和 3 年 swap rate bootstrapping 2 年、3 年 LIBOR zero rates。

\chapter{Securitization / 证券化}
\source{Lecture 9; Assignment 3}

\section{Basic Structure}
\term{Securitization}{证券化} 把一组资产的现金流打包给 SPV，再由 SPV 发行不同层级的证券。典型结构：
\[
\text{Loans/Mortgages}\rightarrow \text{SPV}\rightarrow \text{Senior, Mezzanine, Equity tranches}.
\]

\section{Waterfall}
\term{Waterfall}{瀑布结构} 决定现金流和损失如何分配。通常收益先给 senior，再给 mezzanine，最后给 equity；损失反过来，equity first loss，之后 mezzanine，再之后 senior。

\begin{intuitionbox}{为什么 senior tranche 也可能有风险}
Senior tranche 不是没有风险，而是只有在低层级已经被打穿后才开始亏损。ABS CDO 的危险在于它把很多 ABS mezzanine tranches 再打包，底层相关性一上升，原本看起来分散的 mezzanine 风险会集中冲击 CDO senior tranche。
\end{intuitionbox}

\chapter{Option Properties / 期权性质}
\source{Lecture 10; Assignment 1; Assignment 3}

\section{Call and Put}
\begin{itemize}
  \item \term{Call option}{看涨期权}：持有人有权买入标的。
  \item \term{Put option}{看跌期权}：持有人有权卖出标的。
  \item \term{European option}{欧式期权}：只能到期行权。
  \item \term{American option}{美式期权}：到期前任意时间可行权。
\end{itemize}

\section{Payoff and Profit}
Long call payoff:
\[
\max(S_T-K,0).
\]
Long put payoff:
\[
\max(K-S_T,0).
\]
Profit = payoff - premium accumulated to maturity. 在很多 assignment solution 中，为了简单忽略 time value of money，直接用 payoff - premium。

\section{Bounds}
不支付股利股票的 European call 下界：
\[
c\geq S_0-Ke^{-rT}.
\]
European put 下界：
\[
p\geq Ke^{-rT}-S_0.
\]
同时 $c\leq S_0$，$p\leq Ke^{-rT}$。

\section{Put-Call Parity}
无股利 European options：
\[
c+Ke^{-rT}=p+S_0.
\]
有确定股利现值 $D$：
\[
c+D+Ke^{-rT}=p+S_0.
\]

\begin{intuitionbox}{Parity 的直觉}
左边“call + bond + dividend cash reserve”和右边“put + stock”到期现金流一样。到期如果 $S_T>K$，两边都等于 $S_T$；如果 $S_T<K$，两边都等于 $K$。今天价格不同就有 arbitrage。
\end{intuitionbox}

\section{Early Exercise}
对于不支付股利的股票，American call 通常不应提前行权。原因是行权要提前支付 strike，损失利息，而且仍然没有获得额外股利。American put 可能提前行权，尤其是 deep in the money 时，因为提前拿到 $K$ 可以赚利息。

\chapter{Trading Strategies / 期权交易策略}
\source{Lecture 11; Assignment 3}

\section{Covered Call and Protective Put}
\term{Covered call}{备兑看涨}：long stock + short call。它牺牲上涨空间，换取 premium income。

\term{Protective put}{保护性看跌}：long stock + long put。它像给股票买保险，锁定最低卖出价值。

Put-call parity 可以写成：
\[
S_0+p=c+Ke^{-rT}.
\]
这说明 protective put 的经济效果等价于 fiduciary call（call + risk-free bond）。

\section{Spreads}
\begin{itemize}
  \item \term{Bull spread}{牛市价差}：押温和上涨，收益和亏损都有限。
  \item \term{Bear spread}{熊市价差}：押温和下跌。
  \item \term{Butterfly spread}{蝶式价差}：押价格停留在中间 strike 附近。
  \item \term{Calendar spread}{日历价差}：同 strike 不同 maturity。
\end{itemize}

\section{Combinations}
\begin{itemize}
  \item \term{Straddle}{跨式组合}：同 strike 同 maturity 的 long call + long put，押大波动。
  \item \term{Strip}{条式组合}：one call + two puts，押大波动且更偏向下跌。
  \item \term{Strap}{带式组合}：two calls + one put，押大波动且更偏向上涨。
  \item \term{Strangle}{宽跨式组合}：不同 strike 的 out-of-the-money call 和 put，成本更低但需要更大波动。
\end{itemize}

\chapter{Binomial Trees / 二叉树定价}
\source{Lecture 12; Assignment 4}

\section{One-Step Binomial Model}
假设股票现在为 $S_0$，一个 period 后只可能上涨到 $S_0u$ 或下跌到 $S_0d$。衍生品 payoff 分别为 $f_u$ 和 $f_d$。

复制组合 long $\Delta$ shares and short one derivative，使组合无风险：
\[
S_0u\Delta-f_u=S_0d\Delta-f_d.
\]
因此
\[
\Delta=\frac{f_u-f_d}{S_0u-S_0d}.
\]

\section{Risk-Neutral Valuation}
更常用的写法是 risk-neutral probability：
\[
p=\frac{e^{r\Delta t}-d}{u-d}.
\]
衍生品价值：
\[
f=e^{-r\Delta t}\left[pf_u+(1-p)f_d\right].
\]

\begin{warningbox}{Risk-neutral probability 不是真实概率}
$p$ 不是股票真实上涨概率，而是让 discounted stock price 成为 martingale 的定价权重。考试里不要解释成“市场认为上涨概率是 $p$”。
\end{warningbox}

\section{Two-Step Tree}
两步树可以从终点向前 backward induction：
\[
f_{\text{node}}=e^{-r\Delta t}\left[p f_{\text{up}}+(1-p)f_{\text{down}}\right].
\]
如果是 American option，每个节点还要比较 continuation value 和 immediate exercise value：
\[
f_{\text{American put}}=\max\left(K-S,\; e^{-r\Delta t}[pf_u+(1-p)f_d]\right).
\]

\chapter{Detailed Beginner Notes / 逐章超详细讲义}

这一章是对前面主题章的加厚版。前面章节适合快速复习公式；本章适合从零理解。每一节都按“直觉 - 现金流 - 公式 - 题目入口 - 常见坑”的顺序写。读的时候不要急着背结论，先把每个产品的 long side 和 short side 现金流画出来。

\section{Foundations / 衍生品基础}

\subsection{Derivative 是一张未来现金流规则}
\term{Derivative}{衍生品} 的价值来自 \term{underlying asset}{标的资产}。这句话的意思不是“衍生品价格一定跟标的价格同步涨跌”，而是说衍生品合约规定了未来某些状态下的现金流，而这些状态通常由标的资产价格、利率、汇率或信用事件决定。

例如，一份 call option 给买方未来按 strike price 买入股票的权利。它不是股票本身，但它的 payoff 由未来股票价格 $S_T$ 决定。一份 interest rate swap 也不是债券本身，但它的现金流由未来 floating reference rate 决定。只要你能写出未来 cash flows，就能开始讨论 pricing。

\begin{workflowbox}{小白解题第一步}
\begin{enumerate}
  \item 圈出 underlying asset：股票、汇率、利率、商品还是信用资产。
  \item 圈出 position：long, short, buy, sell, receive, pay。
  \item 写出 future cash inflow 和 cash outflow。
  \item 看能不能用更简单的资产组合复制这些现金流。
\end{enumerate}
\end{workflowbox}

\subsection{No-arbitrage 不是预测，而是相对定价}
这门课的 pricing 不是预测未来价格。No-arbitrage pricing 的逻辑是：如果两个 portfolios 在未来每一种状态下现金流都一样，那么它们今天的价格必须一样。否则可以买便宜的一边、卖贵的一边，锁定 profit。

这就是 forward pricing, put-call parity, swap valuation, binomial tree 的共同核心。公式虽然不同，但背后的句式很像：
\begin{quote}
Portfolio A and Portfolio B have identical future cash flows. Therefore, their current values must be equal under no arbitrage.
\end{quote}

\subsection{Hedger, speculator, arbitrageur 的区别}
\term{Hedger}{套保者} 已经有一个风险，使用 derivatives 是为了 reduce risk。例如航空公司未来要买 fuel，怕 fuel price rises，所以 long futures。农场主未来要卖 corn，怕 corn price falls，所以 short futures。

\term{Speculator}{投机者} 本来没有必须对冲的风险，而是主动承担价格方向风险。看涨就 long， 看跌就 short。Speculation 常用 derivatives 是因为 leverage 高，不需要付出全部 underlying notional。

\term{Arbitrageur}{套利者} 关注 price relation 是否错位。例如 forward price 高于 cost-of-carry value，就可以 cash-and-carry；put-call parity 左边贵，就 short 左边、long 右边。

\begin{warningbox}{基础章最常见误区}
不要把 hedging 写成 eliminate all risk。更准确是 reduce or manage risk，因为 hedge 仍然可能有 basis risk, maturity mismatch, quantity mismatch, liquidity risk。
\end{warningbox}

\section{Futures Markets and Hedging / 期货与套保}

\subsection{Futures 为什么要 daily settlement}
Futures 和 forward 都是在未来按约定价格交易 underlying asset，但 futures 是 exchange traded and standardized，最重要机制是 \term{daily settlement}{每日结算}。每天交易所根据 settlement price 计算当天 gains and losses，并把钱划入或划出 margin account。

\term{Initial margin}{初始保证金} 是开仓时放入账户的钱，不是 purchase price。Futures 合约本身初始价值通常为零。保证金是 performance bond，用来保证你能承担未来 daily losses。\term{Maintenance margin}{维持保证金} 是最低余额；低于这个水平时触发 \term{margin call}{追加保证金}。

\begin{workflowbox}{Margin call 题步骤}
\begin{enumerate}
  \item 计算 total initial margin 和 total maintenance margin。
  \item 两者差额就是触发 margin call 前最多可亏的金额。
  \item 用亏损金额除以 total contract units，得到 price change。
  \item Long futures 怕 price falls；short futures 怕 price rises。
  \item 如果问可 withdraw 多少，账户必须高于 initial margin 对应金额。
\end{enumerate}
\end{workflowbox}

\begin{examplebox}{Assignment 1 Q3}
Long 2 张 FCOJ futures，每张 15,000 pounds。Initial margin 总额是 \$16,000，maintenance margin 总额是 \$12,000，所以亏 \$4,000 触发 margin call。总数量是 30,000 pounds，价格下跌为
\[
\frac{4000}{30000}=0.1333\text{ dollars}=13.33\text{ cents}.
\]
因为是 long futures，价格下跌亏损，所以触发价是 $200-13.33=186.67$ cents/lb。
\end{examplebox}

\subsection{Basis risk 为什么 hedge 后还会不确定}
\term{Basis}{基差} 定义为 spot price minus futures price：
\[
\text{basis}=S-F.
\]
接近 delivery period 时，同一资产的 futures price 和 spot price 应该 converge，否则可以通过交割套利获利。但实际 hedge 往往不是完美的，因为公司真实买卖的资产、地点、质量、时间和期货合约不完全一致。

例如航空公司可能用 heating oil 或 gasoline futures 对冲 jet fuel 风险。这就是 \term{cross hedge}{交叉套保}。Cross hedge 的关键不是完全消灭风险，而是选择相关性高的 futures contract，让 hedge portfolio variance 最小。

\subsection{Minimum variance hedge ratio}
最小方差套保比率：
\[
h^\ast=\rho\frac{\sigma_S}{\sigma_F}.
\]
这个公式的直觉是：
\begin{itemize}
  \item $\rho$ 越高，futures 越能解释 spot risk，hedge ratio 越大。
  \item $\sigma_S$ 越大，被套保资产波动越大，需要更多 futures。
  \item $\sigma_F$ 越大，一张 futures 对风险的抵消能力越强，需要更少 futures。
\end{itemize}

\begin{examplebox}{Assignment 1 Q6}
公司每 1 cent/gallon 的新燃料涨价损失 \$1,000,000，所以 exposure 是
\[
\frac{1000000}{0.01}=100000000\text{ gallons}.
\]
相关系数 $0.8$，新燃料波动率是 gasoline futures 的 $1.5$ 倍，所以 $h^\ast=0.8\times1.5=1.2$。公司怕价格上涨，因此 long futures。需要对冲 $120,000,000$ gallons，每张 40,000 gallons，所以 long 3,000 张。
\end{examplebox}

\subsection{Beta hedge 的符号含义}
用 stock index futures 调整股票组合 beta：
\[
N^\ast=(\beta_T-\beta_P)\frac{V_A}{F_0M}.
\]
如果目标 beta 小于当前 beta，$N^\ast$ 为负，表示 short index futures；如果目标 beta 大于当前 beta，$N^\ast$ 为正，表示 long index futures。这里分母 $F_0M$ 是一张 futures 的 notional exposure，不是保证金。

\section{Interest Rates, Zero Rates and FRA / 利率、零息与 FRA}

\subsection{利率题先看 compounding convention}
同样是 $8\%$，annual compounding, semiannual compounding, continuous compounding 代表不同 future value。衍生品定价里，所有 cash flows 必须用一致的 compounding convention 处理。

若每年复利 $m$ 次：
\[
A=P\left(1+\frac{R_m}{m}\right)^{mT}.
\]
若连续复利：
\[
A=Pe^{R_cT},\qquad P=Ae^{-R_cT}.
\]
离散到连续的转换：
\[
R_c=m\ln\left(1+\frac{R_m}{m}\right).
\]

\begin{warningbox}{利率章常见错误}
不要把 $6$ months 当成 $6$ 代入。所有 $T$ 都用 year fraction，例如 $6$ months $=0.5$，$3$ months $=0.25$。
\end{warningbox}

\subsection{Zero rate 和 discount factor 的关系}
\term{Zero rate}{零息利率} 是今天到某个 maturity 的 single cash flow 利率。若连续复利 zero rate 为 $R(T)$，discount factor 是
\[
P(0,T)=e^{-R(T)T}.
\]
Coupon bond 不是一个 cash flow，而是一串 cash flows，所以价格是
\[
B=\sum_i C_iP(0,t_i).
\]

这也是 bootstrapping 的本质：短端 discount factors 已知后，用更长期债券或 swap 的价格方程，解出下一个未知 discount factor。

\begin{workflowbox}{Bootstrapping 步骤}
\begin{enumerate}
  \item 画时间线，列出每期 coupon 和 final principal。
  \item 已知 maturity 的 cash flow 先折现并扣掉。
  \item 价格方程里只留下一个未知 discount factor。
  \item 解出 discount factor 后，再转换成 zero rate。
\end{enumerate}
\end{workflowbox}

\subsection{Forward rate 不是平均利率}
Forward rate 是今天隐含的未来某段期间利率。设 $R_1$ 是到 $T_1$ 的 zero rate，$R_2$ 是到 $T_2$ 的 zero rate，则连续复利下：
\[
e^{R_2T_2}=e^{R_1T_1}e^{R_F(T_2-T_1)}.
\]
取对数得到：
\[
R_F=\frac{R_2T_2-R_1T_1}{T_2-T_1}.
\]
它不是 $R_1$ 和 $R_2$ 的简单平均。若 yield curve upward sloping，forward rate 可能高于较长期 zero rate。

\subsection{FRA 的结算时点}
Forward Rate Agreement 锁定未来 $T_1$ 到 $T_2$ 的借贷利率。核心 cash flow 是 rate difference times notional times period length。若约定利率为 $R_K$，实际市场利率为 $R_M$，本金 $L$，期间长度 $\tau$，期末利息差是
\[
L(R_M-R_K)\tau
\]
或反方向，取决于你是 receiving fixed 还是 paying fixed。实际 FRA 常在 $T_1$ 结算，因此要把 $T_2$ 的利息差折现回 $T_1$。这就是很多题会多除一个 $1+R_M\tau$ 的原因。

\section{Forward and Futures Pricing / 远期与期货定价}

\subsection{Cash-and-carry 的完整逻辑}
无收益 investment asset 的 forward price：
\[
F_0=S_0e^{rT}.
\]
这不是凭空来的。构造组合 A：今天借钱买入现货并持有到期。到期你拥有一单位 asset，且需要还 $S_0e^{rT}$。组合 B：今天进入 long forward，到期支付 $F_0$ 买入同一单位 asset。两者到期都得到一单位 asset，所以 financing cost 和 forward delivery price 必须一致。

如果市场 $F_0>S_0e^{rT}$，forward 太贵。套利者可以 borrow $S_0$, buy spot, short forward。到期交割资产收到 $F_0$，还贷款 $S_0e^{rT}$，锁定 profit。若 $F_0<S_0e^{rT}$，做 reverse cash-and-carry。

\begin{workflowbox}{Forward pricing 公式选择}
\begin{enumerate}
  \item 没有收益、没有成本：$F_0=S_0e^{rT}$。
  \item 有确定 cash income，现值为 $I$：$F_0=(S_0-I)e^{rT}$。
  \item 有 continuous yield $q$：$F_0=S_0e^{(r-q)T}$。
  \item 外汇：$F_0=S_0e^{(r_d-r_f)T}$。
  \item 商品有 storage cost 现值 $U$：$F_0=(S_0+U)e^{rT}$。
\end{enumerate}
\end{workflowbox}

\subsection{I 和 q 不要混用}
$I$ 是确定现金收益的 present value，例如已知 dividend in dollars。它是金额，直接从 $S_0$ 中减掉。$q$ 是 continuous dividend yield，是年化比例，放在指数里。Stock index futures 常用 $q$，单只股票已知几次现金股利常用 $I$。

\begin{examplebox}{Assignment 2 Q1}
股票 $S_0=100$，2 个月和 5 个月各支付 \$1，$r=8\%$ continuous，6 个月 forward。先算 dividend present value：
\[
I=e^{-0.08(2/12)}+e^{-0.08(5/12)}=1.9540.
\]
再代入
\[
F_0=(100-1.9540)e^{0.08(0.5)}=102.05.
\]
三个月后问的是 old forward contract 的 value，不是重新问初始 forward price，所以要使用旧 delivery price $K$ 并注意 long/short 符号。
\end{examplebox}

\subsection{Currency forward 的 domestic 和 foreign}
外汇远期：
\[
F_0=S_0e^{(r_d-r_f)T}.
\]
这里 $S_0$ 的报价方向决定 domestic 和 foreign。若 $S_0$ 是 USD per EUR，则 domestic currency 是 USD，foreign currency 是 EUR。Foreign currency 可以看成一项提供 foreign risk-free yield 的资产，所以 $r_f$ 像 dividend yield 一样降低 forward price。

\subsection{Commodity: storage cost 和 convenience yield}
商品定价比股票更复杂，因为商品可能有 storage cost，也可能有 convenience yield。Storage cost 提高持有现货成本，因此提高 futures price。Convenience yield 是持有实物带来的非现金收益，例如保证生产不断、避免短缺，所以降低 futures price。

连续形式：
\[
F_0=S_0e^{(r+u-y)T}.
\]
对 consumption asset，这个公式经常不是 exact price，而是 upper bound。因为 reverse cash-and-carry 可能需要 short sell commodity，但实物商品并不总能轻易卖空或借入。

\subsection{Contango 和 backwardation}
\term{Contango}{正向市场} 通常指远月价格高于近月价格，常见于 financing cost 和 storage cost 占主导时。\term{Backwardation}{反向市场} 指近月价格高于远月价格，常见于短期供给紧张或 convenience yield 很高时。

Oil futures supplement 的重点不是记某一天油价，而是理解 futures curve 反映库存、运输、地缘风险和短期 scarcity。地缘冲击通常先影响 near-month contracts，因为市场最担心眼前 supply disruption。

\section{Interest Rate Futures / 利率期货}

\subsection{Bond quoted price 和 cash price}
Treasury bond 的 quoted price 通常不含 accrued interest。买方真正支付的是
\[
\text{Cash price}=\text{Quoted price}+\text{Accrued interest}.
\]
报价 102-07 表示 $102+7/32=102.21875$，不是 $102.07$。Accrued interest 取决于 coupon period 中已经过去的天数。

\begin{workflowbox}{债券报价题步骤}
\begin{enumerate}
  \item 把 32nds quote 转成 decimal quote。
  \item 找上一付息日到 settlement date 的 elapsed days。
  \item 计算 accrued interest = coupon payment times elapsed days / coupon period days。
  \item Cash price = quoted price + accrued interest。
\end{enumerate}
\end{workflowbox}

\subsection{Conversion factor 和 CTD}
Treasury bond futures 允许 short party 选择一篮子 deliverable bonds 中的一只交割。为了让不同 coupon 和 maturity 的债券可比，交易所有 conversion factor。

交割时 short 收到的 invoice amount 近似为：
\[
\text{Invoice amount}=\text{Futures settlement price}\times\text{conversion factor}+\text{accrued interest}.
\]
Short 会选 \term{cheapest-to-deliver}{CTD} bond，即交割净成本最低的债券。CTD 不是价格最低的债券，而是 relative delivery economics 最优的债券。

\subsection{Eurodollar futures 的方向}
Eurodollar futures quote:
\[
Q=100-\text{implied 3-month LIBOR}.
\]
所以 quote 上升表示利率下降，quote 下降表示利率上升。借款人怕未来利率上升，利率上升时期货价格下降，因此借款人应该 short Eurodollar futures。

\begin{examplebox}{Assignment 2 Q8}
Quote $98.40$ 表示 implied LIBOR 为 $1.60\%$。公司未来按 LIBOR + $0.5\%$ 借款，所以锁定 borrowing rate 约 $2.10\%$。如果实际 LIBOR 是 $1.3\%$，最终 settlement price 是 $98.70$。
\end{examplebox}

\subsection{Duration hedge}
债券价格对 yield 的敏感度近似为：
\[
\frac{\Delta B}{B}\approx-D\Delta y.
\]
债券组合用 Treasury bond futures 对冲时：
\[
N^\ast=\frac{PD_P}{FD_F}.
\]
如果你持有 long bond portfolio，担心 rates rise 导致 bond prices fall，通常 short futures。Rates rise 时 futures price 也下跌，short futures gain 抵消债券损失。

\section{Swaps / 互换}

\subsection{Plain vanilla swap 的现金流}
Plain vanilla interest rate swap 是一方 pay fixed and receive floating，另一方 receive fixed and pay floating。Notional principal 通常不交换，只用于计算利息。

如果公司原本有 floating-rate loan，比如支付 LIBOR + $0.8\%$，同时进入 pay fixed $5\%$ receive LIBOR 的 swap，那么 LIBOR 支出和 LIBOR 收入抵消，最终近似固定成本为 $5.8\%$。

\subsection{Swap 可以拆成 bonds}
对 pay fixed, receive floating 的一方：
\[
V_{\text{swap}}=B_{\text{floating}}-B_{\text{fixed}}.
\]
直觉是：你收 floating cash flows，好像 long 一只 floating-rate bond；你付 fixed cash flows，好像 short 一只 fixed-rate bond。Receive fixed, pay floating 的一方符号相反。

Floating-rate bond 在 reset date 附近通常约等于 par，因为下一期 coupon 会按市场利率重置。Fixed-rate bond 要把所有固定 coupon 和 principal 用当前 zero curve 折现。

\subsection{Swap 也可以拆成一串 FRA}
每一期 swap payment 都像一个 FRA：固定利率和未来浮动利率之间的差额乘 notional 再乘 year fraction。把所有未来净现金流折现求和，就得到 swap value。

Par swap rate 是让初始 swap value 为零的 fixed rate。若 discount factors 已知，支付间隔为 $\alpha_i$，常见形式是
\[
R_{\text{swap}}=\frac{1-P(0,T_n)}{\sum_{i=1}^{n}\alpha_iP(0,T_i)}.
\]
Assignment 3 Q1 用 par swap 等价 par coupon bond 的方法，从 2 年和 3 年 swap rates 推出 zero rates。

\subsection{Comparative advantage 题}
比较优势题先比较两家公司在 fixed market 和 floating market 的 borrowing spread。两个 spread 的差额是 total potential gain。然后让每家公司在相对有优势的市场借款，再通过 swap 达到各自真正想要的 fixed 或 floating exposure。

\begin{workflowbox}{Comparative advantage 题步骤}
\begin{enumerate}
  \item 计算 fixed borrowing spread。
  \item 计算 floating borrowing spread。
  \item Quality spread differential = 两个 spread 的差。
  \item 扣除 financial intermediary fee。
  \item 分配剩余 gain，并检查双方最终成本都优于直接借款。
\end{enumerate}
\end{workflowbox}

\section{Securitization / 证券化}

\subsection{证券化不是消灭风险，而是重新切分风险}
Securitization 把 loans 或 mortgages 的现金流转给 SPV，再由 SPV 发行不同 seniority 的 securities。SPV 的作用是隔离 originator 的破产风险，让投资者主要面对 collateral pool 的现金流风险。

Tranches 是风险分层。Senior tranche 优先拿现金流，风险较低；equity tranche 最后拿剩余收益，但承担 first loss。Mezzanine tranche 位于中间。

\subsection{Waterfall 的损失方向}
正常现金流通常从 senior 往下分配；损失从 equity 往上吸收。做损失题时，不能把 total mortgage loss 直接乘到每个 tranche 上。必须先打穿 equity，再打 mezzanine，再打 senior。

\begin{examplebox}{Assignment 3 Q5 第一层 ABS}
ABS: senior $75\%$, mezzanine $20\%$, equity $5\%$。Mortgage loss $16\%$。
\begin{itemize}
  \item Equity thickness $5\%$，先全部损失，所以 equity loss rate $100\%$。
  \item 剩余 loss $16\%-5\%=11\%$ 进入 mezzanine。
  \item Mezzanine principal $20\%$，loss rate $11/20=55\%$。
  \item Senior 还没有被打到，loss rate $0\%$。
\end{itemize}
\end{examplebox}

\subsection{ABS CDO 的风险放大}
ABS CDO 的 collateral 不是原始 mortgages，而是 ABS 的 mezzanine tranches。第一层 ABS mezzanine 的损失率会成为 CDO collateral loss。Assignment 3 Q5 中，ABS mezzanine loss 是 $55\%$，所以 CDO collateral loss 是 $55\%$。CDO equity $5\%$ 和 mezzanine $20\%$ 被完全打穿，CDO senior 承担
\[
\frac{55-25}{75}=40\%
\]
的损失。这说明再证券化会集中 mezzanine risk，并让看似 senior 的 tranche 也暴露于严重损失。

\section{Option Properties / 期权性质}

\subsection{Payoff, profit, exercise condition}
Long call payoff:
\[
\max(S_T-K,0).
\]
Long put payoff:
\[
\max(K-S_T,0).
\]
Payoff 是到期合约支付；profit 是 payoff 减 premium 的成本。如果考虑 time value，premium 还要累积到 maturity。简单图形题常直接 payoff minus premium。

Exercise condition 和 profit condition 不一样。Call 只要 $S_T>K$ 就会 exercise，但买方真正盈利需要 $S_T>K+\text{premium}$。Put 只要 $S_T<K$ 就会 exercise，但盈利需要 $S_T<K-\text{premium}$。

\subsection{Option bounds}
European call on non-dividend stock 的下界：
\[
c\geq \max(S_0-Ke^{-rT},0).
\]
European put 的下界：
\[
p\geq \max(Ke^{-rT}-S_0,0).
\]
若价格低于下界，就可以买入 underpriced option 并卖出复制组合进行 arbitrage。

\subsection{Put-call parity 的复制逻辑}
无股利 European options:
\[
c+Ke^{-rT}=p+S_0.
\]
有确定股利现值 $D$:
\[
c+D+Ke^{-rT}=p+S_0.
\]
左边是 call + cash for strike + dividend reserve，右边是 put + stock。到期无论 $S_T>K$ 还是 $S_T<K$，两边 cash flow 相同。因此今天价格必须相同。

\begin{examplebox}{Assignment 3 Q8}
Call 和 put 都为 \$3，$K=20$，$T=0.25$，$r=10\%$，$S_0=19$，一个月后有 \$1 dividend。左边
\[
3+e^{-0.10/12}+20e^{-0.10(0.25)}\approx23.50,
\]
右边 $3+19=22$。左边贵，所以 sell left side and buy right side：short call, borrow present value of strike, buy put and stock。
\end{examplebox}

\subsection{American option 的提前行权}
American option 的 value 至少等于 European option，因为它多了提前行权权利。不支付股利股票的 American call 通常不提前行权，因为提前行权会失去 time value，并提前支付 strike。American put 可能提前行权，因为深度实值时提前拿到 $K-S$ 可以赚利息。

\section{Trading Strategies / 期权交易策略}

\subsection{策略题不要背图，先逐项相加}
Trading strategies 的名字很多，但计算方法很统一：列出每个 option leg，写 payoff，按 strikes 把 $S_T$ 分区间，然后加总。

\begin{workflowbox}{Option strategy 题步骤}
\begin{enumerate}
  \item 列出每个 leg：long/short, call/put, strike, premium。
  \item 写出每个 leg 的 payoff，short option 要取负。
  \item 按所有 strikes 切分 $S_T$ 区间。
  \item 每个区间加总 payoff。
  \item 加上收到的 premium 或减去支付的 premium，得到 profit。
\end{enumerate}
\end{workflowbox}

\subsection{Covered call 和 protective put}
Covered call = long stock + short call。它适合持有股票但认为上涨空间有限的投资者。卖 call 收 premium，但股票涨过 strike 后，上方收益被 short call 抵消。

Protective put = long stock + long put。它像保险。如果股票大跌，put payoff 提供最低卖出保护；如果股票上涨，仍保留 upside，但成本是 premium。

\subsection{Spreads 和 combinations}
Bull spread 通常 buy low strike call and sell high strike call，押温和上涨。Bear spread 押温和下跌。Butterfly spread 通常 buy low strike, buy high strike, sell two middle strike options，押到期价格靠近中间 strike。

Straddle = long call + long put, same strike and same maturity，押大波动不押方向。Strangle 也是押大波动，但 call 和 put 的 strikes 不同，成本通常更低，要求价格移动更大。Short straddle 或 short strangle 是押低波动，但极端行情下亏损可能很大。

\begin{examplebox}{Assignment 3 Q9}
Put strikes 55, 60, 65，价格 4, 6, 9。构造 butterfly：buy 55 put, buy 65 put, sell two 60 puts。净成本
\[
4+9-2(6)=1.
\]
Profit 在 $S_T=60$ 附近最大，亏损区间是 $S_T<56$ 或 $S_T>64$。
\end{examplebox}

\section{Binomial Trees / 二叉树}

\subsection{一步树的两种理解}
Binomial tree 假设一个 period 后股票只会上涨到 $S_0u$ 或下跌到 $S_0d$。定价有两种等价方法：replicating portfolio 和 risk-neutral valuation。

复制组合方法找 $\Delta$ shares 和 cash position，使组合在 up/down 两个状态都等于 derivative payoff。由于未来完全复制成功，今天价格必须等于复制组合今天的成本。

Risk-neutral valuation 是更快的写法：
\[
p=\frac{e^{r\Delta t}-d}{u-d},\qquad
f=e^{-r\Delta t}[pf_u+(1-p)f_d].
\]
$p$ 不是真实上涨概率，而是 pricing weight。

\subsection{两步树从终点往前推}
多步树先画 stock price tree，再在 maturity 写 terminal payoff，然后 backward induction。每一个中间节点都像一个一步树：
\[
f_{\text{node}}=e^{-r\Delta t}[pf_u+(1-p)f_d].
\]
如果是 recombining tree，先上后下和先下后上到达同一节点，可以减少计算。

\begin{examplebox}{Assignment 4 Q1}
$S_0=50$, $u=1.06$, $d=0.95$, $K=51$, $\Delta t=0.25$, $r=5\%$。先算
\[
p=\frac{e^{0.05(0.25)}-0.95}{1.06-0.95}=0.568895.
\]
终点 stock prices 是 $56.18$, $50.35$, $45.125$。Call payoffs 是 $5.18,0,0$，向前折现得到 call value 约 $1.64$。
\end{examplebox}

\subsection{American option 每个节点都要比}
American option 不只是到初始节点多比较一次，而是每个可行权节点都要比较：
\[
\text{American value}=\max(\text{immediate exercise value},\text{continuation value}).
\]
American put 的 immediate exercise value 是 $K-S$。Assignment 4 Q2 中 down node 的股价是 $47.5$，immediate exercise 是 $51-47.5=3.5$，continuation value 约 $2.87$，所以这个节点应该提前行权。

\begin{warningbox}{Assignment 4 复习重点}
Assignment 4 没有官方 solution PDF。本讲义中的解析按 Lecture 12 的方法推导。考试复习时不要只背答案，要能独立完成四步：画 stock tree，写 terminal payoff，算 risk-neutral probability，backward induction。American put 还要在每个节点比较 early exercise。
\end{warningbox}


\chapter{Formula Sheet / 公式速查}
\begin{longtable}{p{0.28\textwidth}p{0.32\textwidth}p{0.32\textwidth}}
\toprule
\textbf{Topic} & \textbf{Formula} & \textbf{Use} \\
\midrule
Continuous compounding & $A=Pe^{RT}$ & discounting and future value \\
Discrete to continuous & $R_c=m\ln(1+R_m/m)$ & rate conversion \\
Forward rate & $R_F=(R_2T_2-R_1T_1)/(T_2-T_1)$ & implied future rate \\
No income forward & $F_0=S_0e^{rT}$ & investment asset with no income \\
Known income forward & $F_0=(S_0-I)e^{rT}$ & stock/bond with known cash income \\
Dividend yield & $F_0=S_0e^{(r-q)T}$ & index futures \\
Currency forward & $F_0=S_0e^{(r_d-r_f)T}$ & FX forward \\
Storage cost & $F_0=(S_0+U)e^{rT}$ & commodity upper bound \\
Convenience yield & $F_0=S_0e^{(r+u-y)T}$ & commodity futures \\
Minimum variance hedge & $h^\ast=\rho\sigma_S/\sigma_F$ & cross hedge \\
Stock index hedge & $N^\ast=(\beta_T-\beta_P)V_A/(F_0M)$ & beta adjustment \\
Duration hedge & $N^\ast=PD_P/(FD_F)$ & bond portfolio hedge \\
Put-call parity & $c+D+Ke^{-rT}=p+S_0$ & European option arbitrage \\
Binomial probability & $p=(e^{r\Delta t}-d)/(u-d)$ & risk-neutral pricing \\
Binomial value & $f=e^{-r\Delta t}[pf_u+(1-p)f_d]$ & option/tree valuation \\
\bottomrule
\end{longtable}

\chapter{High-Frequency Practice Questions / 高频重要题}

\section{Assignment 1 Q1: Short Forward on GBP}
\textbf{Question.} An investor enters into a short forward contract to sell 100,000 British pounds at 1.4000 USD/GBP. Gain or loss if final exchange rate is 1.3800 or 1.4300?

\textbf{Solution.} Short forward means you must sell GBP at 1.4000. If market is 1.3800, selling at 1.4000 is better:
\[
(1.4000-1.3800)\times100000=2000.
\]
Gain = \$2,000. If market is 1.4300, selling at 1.4000 is worse:
\[
(1.4300-1.4000)\times100000=3000.
\]
Loss = \$3,000.

\section{Assignment 1 Q3: Margin Call}
\textbf{Question.} A trader buys two July futures contracts on frozen orange juice. Each contract is 15,000 pounds. Current futures price is 200 cents/lb, initial margin \$8,000 per contract, maintenance margin \$6,000 per contract. What price change leads to margin call? When can \$3,000 be withdrawn?

\textbf{Solution.} Total initial margin $=16000$, total maintenance margin $=12000$. Loss that triggers margin call is \$4,000. Two contracts cover $30000$ pounds, so price fall:
\[
\frac{4000}{30000}=0.1333\text{ dollars}=13.33\text{ cents}.
\]
Margin call when price falls to $200-13.33=186.67$ cents/lb. To withdraw \$3,000, account must rise to \$19,000, so gain is \$3,000:
\[
\frac{3000}{30000}=10\text{ cents}.
\]
Price must rise to 210 cents/lb.

\section{Assignment 1 Q6: Cross Hedge}
\textbf{Question.} A company loses \$1 million for each 1 cent/gallon increase in a new fuel. Correlation with gasoline futures price changes is 0.8. New fuel standard deviation is 50\% greater than gasoline futures. Contract size is 40,000 gallons.

\textbf{Solution.} Minimum variance hedge ratio:
\[
h^\ast=0.8\times1.5=1.2.
\]
Exposure measured in gallons:
\[
\frac{1000000}{0.01}=100000000\text{ gallons}.
\]
Because the company loses when fuel price rises, it should gain when gasoline futures rises, so take a long futures position:
\[
1.2\times100000000=120000000\text{ gallons}.
\]
Number of contracts:
\[
\frac{120000000}{40000}=3000.
\]

\section{Assignment 1 Q7: Changing Beta}
Portfolio value is \$100 million, beta is 1.4, index is 2,000, multiplier is \$250.

\textbf{To reduce beta to 0.5:}
\[
N=(0.5-1.4)\frac{100000000}{2000\times250}=-180.
\]
Short 180 futures contracts.

\textbf{To increase beta to 1.8:}
\[
N=(1.8-1.4)\frac{100000000}{2000\times250}=80.
\]
Long 80 futures contracts.

\section{Assignment 2 Q1: Forward Price with Dividends}
Stock price $S_0=100$, dividends \$1 in 2 months and 5 months, $r=8\%$, $T=0.5$.

Present value of dividends:
\[
I=e^{-0.08(2/12)}+e^{-0.08(5/12)}=1.9540.
\]
Forward price:
\[
F_0=(100-1.9540)e^{0.08(0.5)}=102.05.
\]
Initial value of the forward is zero.

Three months later, $S=48$, one dividend remains in 2 months:
\[
I=e^{-0.08(2/12)}=0.9868.
\]
The short forward value is
\[
f_{\text{short}}=-\left(S-I-Ke^{-rT}\right)=51.05.
\]
New 3-month forward price:
\[
F=(48-0.9868)e^{0.08(3/12)}=47.96.
\]

\section{Assignment 2 Q2: Stock Index Futures}
If $S_0=1200$, $r=9\%$, average dividend yield $q=3.4\%$, and $T=153/365$:
\[
F_0=1200e^{(0.09-0.034)(153/365)}\approx1228.5.
\]

\section{Assignment 2 Q3: Currency Forward}
Given $S_0=1.4500$, $F_0=1.3950$, USD rate $r_d=1\%$, $T=0.5$:
\[
1.3950=1.4500e^{(0.01-r_f)0.5}.
\]
So
\[
r_f=0.01-\frac{\ln(1.3950/1.4500)}{0.5}\approx8.73\%.
\]

\section{Assignment 2 Q4-Q5: Commodity Futures}
Oil: $S_0=80$, storage cost \$3 paid at the end of one year, $r=5\%$:
\[
F_0\leq80e^{0.05}+3=87.10.
\]
It is an upper bound because oil is a consumption asset and convenience yield may reduce futures price.

Silver: $S_0=15$, storage \$0.24 per year paid quarterly in advance, $T=0.75$, $r=10\%$:
\[
U=0.06+0.06e^{-0.10(0.25)}+0.06e^{-0.10(0.5)}=0.1756.
\]
\[
F_0=(15+0.1756)e^{0.10(0.75)}\approx16.36.
\]

\section{Assignment 2 Q8: Eurodollar Futures Hedge}
December Eurodollar futures quote is 98.40, so implied LIBOR is $100-98.40=1.60\%$. Borrowing at LIBOR + 0.5\% locks in about $2.10\%$.

A borrower should short Eurodollar futures. If actual 3-month rate is 1.3\%, final settlement price is
\[
100-1.3=98.70.
\]

\section{Assignment 3 Q1: Swap Rates to Zero Rates}
One-year LIBOR zero rate is 10\% annually compounded. Two-year swap rate is 11\%. Treat the par swap as a par coupon bond:
\[
\frac{11}{1.10}+\frac{111}{(1+R_2)^2}=100.
\]
Thus $R_2=11.05\%$.

For the 3-year swap rate 12\%:
\[
\frac{12}{1.10}+\frac{12}{(1.1105)^2}+\frac{112}{(1+R_3)^3}=100.
\]
Thus $R_3=12.17\%$.

\section{Assignment 3 Q5: ABS CDO Loss Waterfall}
Mortgage portfolio loss is 16\%. ABS has equity 5\%, mezzanine 20\%, senior 75\%.

ABS losses:
\begin{itemize}
  \item Equity tranche: 100\%.
  \item Mezzanine tranche: $(16-5)/20=55\%$.
  \item Senior tranche: 0\%.
\end{itemize}

ABS CDO is created from mezzanine tranches, so its collateral loss is 55\%. With CDO equity 5\%, mezzanine 20\%, senior 75\%:
\begin{itemize}
  \item CDO equity: 100\%.
  \item CDO mezzanine: 100\%.
  \item CDO senior: $(55-25)/75=40\%$.
\end{itemize}

\section{Assignment 3 Q6-Q7: Option Profit Conditions}
Long European put: premium \$3, strike \$40. Exercise if $S_T<40$. Profit if
\[
40-S_T-3>0 \Rightarrow S_T<37.
\]

Short European call: premium \$4, strike \$50. The option holder exercises if $S_T>50$. Seller profit is
\[
4-\max(S_T-50,0).
\]
Seller profits if $S_T<54$.

\section{Assignment 3 Q8: Put-Call Parity Arbitrage}
Call and put both cost \$3, $K=20$, $T=0.25$, $r=10\%$, $S_0=19$, dividend \$1 in one month.

Parity says
\[
c+D+Ke^{-rT}=p+S_0.
\]
Left side:
\[
3+e^{-0.10/12}+20e^{-0.10(0.25)}\approx23.50.
\]
Right side:
\[
3+19=22.
\]
Left side is overpriced. Arbitrage: sell call, borrow present value of strike, buy put and stock. The future stock/dividend cash flows repay the synthetic obligations and lock in the mispricing.

\section{Assignment 3 Q9-Q10: Butterfly Spread}
Put prices for strikes 55, 60, 65 are 4, 6, 9. Buy one 55 put, buy one 65 put, sell two 60 puts. Initial cost:
\[
4+9-2(6)=1.
\]
Net profit is:
\[
\begin{cases}
-1, & S_T\leq55,\\
S_T-56, & 55<S_T\leq60,\\
64-S_T, & 60<S_T\leq65,\\
-1, & S_T>65.
\end{cases}
\]
Loss occurs when $S_T<56$ or $S_T>64$. Put-call parity shows the cost is identical to the same butterfly built with calls because the bond and stock terms cancel when combining strikes symmetrically.

\section{Assignment 4 Q1-Q2: Two-Step Binomial Tree}
This assignment has no official solution PDF in the materials. The following is a non-official solution derived from Lecture 12.

Given $S_0=50$, $u=1.06$, $d=0.95$, $\Delta t=0.25$, $r=5\%$, $K=51$:
\[
p=\frac{e^{0.05(0.25)}-0.95}{1.06-0.95}=0.568895.
\]
Terminal stock prices:
\[
S_{uu}=56.18,\qquad S_{ud}=50.35,\qquad S_{dd}=45.125.
\]
Call payoffs:
\[
5.18,\quad0,\quad0.
\]
European call value:
\[
c=e^{-0.05(0.5)}\left[p^2(5.18)\right]\approx1.64.
\]

Put payoffs:
\[
0,\quad0.65,\quad5.875.
\]
European put:
\[
p_{\text{put}}\approx1.38.
\]
Put-call parity check:
\[
c+Ke^{-0.05(0.5)}=1.635+49.741\approx51.376,
\]
\[
p+S_0=1.376+50=51.376.
\]
For American put, compare continuation and early exercise. At the down node, immediate exercise is $51-47.5=3.5$, continuation value is about 2.87, so early exercise is optimal there. American put value is about 1.65.

\section{\texorpdfstring{Assignment 4 Q3: One-Step Derivative Paying $S_T^2$}{Assignment 4 Q3: One-Step Derivative Paying ST squared}}
This problem uses $S_0=25$, $S_u=27$, $S_d=23$, $r=10\%$, $T=2/12$. Here $u=27/25$, $d=23/25$:
\[
p=\frac{e^{0.10(2/12)}-23/25}{27/25-23/25}=0.60504.
\]
Payoff is $S_T^2$, so
\[
f_u=27^2=729,\qquad f_d=23^2=529.
\]
Value:
\[
f=e^{-0.10(2/12)}\left[0.60504(729)+(1-0.60504)(529)\right]\approx639.26.
\]

\backmatter
\chapter*{Final Exam Checklist / 考前检查}
\addcontentsline{toc}{chapter}{Final Exam Checklist / 考前检查}
\begin{multicols}{2}
\begin{itemize}
  \item Can you explain no-arbitrage without memorizing?
  \item Can you convert continuous and semiannual rates?
  \item Can you bootstrap a zero rate from bond/swap prices?
  \item Can you price a forward with income, dividend yield, storage cost, or FX yield?
  \item Can you choose long vs short futures for a hedge?
  \item Can you calculate margin call price levels?
  \item Can you compute stock index beta hedge contracts?
  \item Can you compute Eurodollar implied rate from quote?
  \item Can you value a swap as fixed bond vs floating bond?
  \item Can you allocate losses through ABS/CDO tranches?
  \item Can you write put-call parity with dividends?
  \item Can you draw and interpret option strategy payoffs?
  \item Can you run a binomial tree backward?
  \item Can you decide American put early exercise?
\end{itemize}
\end{multicols}

\end{document}
